\begin{helper}
Let $k$ be a natural number and let $w_{ij}\ge 0$ be a nonnegative real
weight assigned to each ordered pair of indices in $\{1,\ldots,k\}$.  For an
increasing triple $i<j<\ell$, put
\[
T(i,j,\ell)=\sqrt{w_{ij}w_{i\ell}w_{j\ell}}.
\]
The ordered triangle sum obtained by taking, for each ordered triple of
distinct indices, the same value $T(a,b,c)$ after sorting its entries
$a<b<c$, is exactly six times the increasing-triple sum:
\[
\begin{aligned}
&\sum_{i<j<\ell}T(i,j,\ell)
+\sum_{i<\ell<j}T(i,\ell,j)
+\sum_{j<i<\ell}T(j,i,\ell)\\
&\quad
+\sum_{j<\ell<i}T(j,\ell,i)
+\sum_{\ell<i<j}T(\ell,i,j)
+\sum_{\ell<j<i}T(\ell,j,i)
=
6\sum_{i<j<\ell}T(i,j,\ell).
\end{aligned}
\]
\end{helper}
\begin{proof}
The nonnegativity hypothesis ensures that these square-root weights are in
the paper-supported setting of nonnegative edge weights.  The identity itself
is a finite reindexing.  Each of the six sums on the left is indexed by one of the six possible
linear orderings of an ordered triple with three distinct entries.  The
first is already the increasing-triple sum.  For the second ordering
$i<\ell<j$, the coordinate swap $(i,j,\ell)\mapsto(i,\ell,j)$ is a bijection
from those triples to the increasing triples, and the summand becomes
$T(i,\ell,j)$ with its arguments increasing.  The same coordinate
permutation gives bijections from the remaining four orderings
$j<i<\ell$, $j<\ell<i$, $\ell<i<j$, and $\ell<j<i$ to the increasing triples.
In every case the displayed summand is exactly the value attached to the
sorted triple.  Hence all six summands are equal to
$\sum_{i<j<\ell}T(i,j,\ell)$, and adding the six equal contributions gives
$6\sum_{i<j<\ell}T(i,j,\ell)$.
\end{proof}
